% =========================================================
\section{引言}
\label{sec:introduction}
% =========================================================

随着大规模储能系统逐渐承担调峰、调频和新能源消纳等关键任务，
储能电站的运行维护正在从传统的人工巡检逐步走向智能诊断与辅助决策。
与开放域问答不同，运维任务并不是从一份静态文档中寻找一个局部答案：
一次故障处置往往需要同时参考设备状态、告警与巡检记录、历史工单、
操作规程以及后续验证结果，并根据这些证据之间的先后关系判断故障如何演化、
哪些判断已经被后续记录修正，以及当前应当执行何种操作。
因此，对于面向储能电站具身运维的智能系统而言，
关键问题不仅是``检索哪些证据与查询语义相关''，
还包括``在当前查询时刻哪些证据已经可知''以及
``这些证据能否共同组成支持当前决策的完整过程''。
系统的最终任务也不是生成一段自由文本答案，
而是将可信证据进一步组织为一个字段完整、格式规范的标准化工单，
以及一个满足规程、安全依赖和物理执行顺序的有序具身指令序列。

检索增强生成（Retrieval-Augmented Generation, RAG）
通过在大语言模型生成前检索外部知识，
有效缓解了模型参数知识过时以及领域知识不足的问题
\cite{lewis2020rag}。
现有 RAG 系统通常首先根据查询与文档之间的语义或词法相关性获得候选，
再将排名靠前的若干文本片段作为生成上下文。
这一范式适用于大量静态知识问答任务，
但直接应用于持续更新的运维记录时存在两个明显局限。

首先，运维知识具有显著的\emph{双时间属性}。
一条记录既有其描述事件实际发生的时间，
也有其进入信息系统并能够被检索到的时间。
晚到的巡检报告、补录工单以及事后形成的诊断结论，
可能描述较早发生的事件，却不应被更早时刻的查询提前使用。
同时，同一设备状态和故障认识还可能随着后续检查不断更新甚至被推翻。
因此，仅按照语义相关性或简单的新鲜度排序，
难以保证检索结果符合查询时刻真实可获得的知识状态。

其次，一次运维决策通常依赖多个相互关联的证据。
例如，异常状态可以触发检查，检查结果支持某一诊断，
诊断进一步引出维修动作，而维修后的复测记录又用于验证处置结果。
这些证据单独来看都可能不是查询中最相似的文本，
但组合起来却构成完成任务所必需的证据链。
传统 Top-$k$ 检索通常独立评价每条文档，
容易得到若干局部相关但功能重复的记录，
同时遗漏连接观察、诊断、处置和验证过程的关键证据。
更一般地说，目标并不一定是一条严格线性的链，
而可能是包含分支、更新、限定与不确定性保留的有向证据结构。

近年来，研究者开始显式研究 RAG 中的时间因素。
MRAG 将查询拆分为语义内容与时间约束，
分别计算证据的语义相关性和时间相关性，
以改善时间敏感问题的检索效果
\cite{zhang2025mrag}；
TA-RAG 则进一步面向跨时间区间的 diachronic question answering，
联合主题语义与时间窗口寻找时间上连贯的证据集合
\cite{lau2025tarag}。
这些研究表明，将时间条件显式引入检索过程
比仅依赖通用语义相似度更适合动态知识环境。
另一方面，多跳检索和结构化 RAG 也说明，
复杂问题往往需要联合多个互补证据，而不是独立选择若干相似片段。
Multi-Hop Dense Retrieval 通过逐跳条件化检索为复杂问题收集多段证据
\cite{xiong2021mdr}；
RAPTOR 通过递归聚类与摘要构建层次化检索结构
\cite{sarthi2024raptor}；
GraphRAG 则利用实体图与社区摘要组织大规模文本语料
\cite{edge2024graphrag}。
然而，这些结构主要来自文档层级、实体联系或多跳问答过程，
并不直接表示持续更新的运维记录之间
\emph{支持、更新、限定、冲突、替代和验证}等查询相关关系。
当一个查询需要恢复从状态变化到处置结果的完整过程时，
时间相关的单条证据排序或预先构建的静态结构仍不能直接保证所选证据集合在流程上完整。

针对上述问题，本文提出
\textbf{TEF-RAG（Temporal Evidence Flow Retrieval-Augmented Generation）}，
面向储能电站运维记录构建查询条件化的时序证据流。
TEF-RAG 不再将候选记录视为彼此独立的文档，
而是首先在查询时刻可见的候选证据之间建立有向语义关系，
并将支持、更新、限定、验证、替代和不确定性传播等关系组织为
\emph{Temporal Evidence Flow}。
在此基础上，系统从候选证据中联合选择一组能够覆盖任务所需角色、
保持时间方向并形成有效关系结构的记录，
而非单纯截取相关性最高的若干文档。

具体而言，TEF-RAG 首先通过基础检索得到候选证据，
并保留资产范围、事件时间与可用时间、
规程适用性等确定性条件作为候选合法性的基本边界。
这些规则主要用于排除在当前查询下不可使用的证据。
随后，系统利用查询条件化的关系判断识别候选证据之间的语义联系，
并通过学习式 pair proposal 减少需要进行关系判断的证据对数量。
最后，与传统逐条文档打分不同，
TEF-RAG 对完整候选证据集合进行联合评价：
采用基于 pairwise ranking 的非线性集合效用模型，
结合困难负样本挖掘，
学习区分能够形成完整证据流的集合与
语义上相关但结构不完整的集合，
最终选择固定规模的证据流，并将其作为下游生成器的受控上下文，
生成标准化工单和有序具身指令序列。
需要说明的是，本文的核心算法贡献仍集中于这一端到端任务中的
\emph{检索与证据组织层}：
检索实验检验系统能否在正确的时间边界内恢复完整、连贯的证据流；
在此基础上，本文进一步固定同一下游生成器，对五种检索上下文进行结构化工单与动作计划评价，
用于检验上游 evidence-flow 质量是否能够传递为更好的生成质量。
该下游评价不等同于真实机器人执行有效性或闭环安全验证。

本文的主要贡献如下：

\begin{enumerate}
    \item
    \textbf{提出面向动态运维知识的时序证据流检索问题。}
    区别于独立文档相关性排序，
    本文从查询时刻可见性和证据间有向关系两个层面描述动态运维知识，
    将检索目标从单条证据相关性扩展为证据集合的流程完整性。

    \item
    \textbf{提出查询条件化的证据关系建模与候选关系生成机制。}
    在基础候选检索之后，
    TEF-RAG 对证据之间的支持、更新、限定、验证、替代等关系进行语义判断，
    并利用学习式 pair proposal 在受控计算预算下
    优先选择与当前查询相关的候选证据对，
    从而构建稀疏的时序证据流。

    \item
    \textbf{提出面向完整证据集合的非线性排序方法。}
    本文不独立预测每条证据的重要性，
    而是在固定候选集合空间中直接学习不同证据组合之间的相对效用，
    并通过 query-level pairwise ranking 与困难负样本挖掘
    建模证据角色、时间结构和关系结构之间的非线性交互，
    使最终选择目标与完整证据流检索更加一致。

    \item
    \textbf{构建并封存面向时序运维证据检索的受控评价协议。}
    评价同时考察传统的 Recall 和 nDCG，
    以及证据组完整性、时序关系完整性等集合级指标，
    并严格区分开发、验证与封存测试阶段。
    该设置用于分析普通相关性检索、时间感知检索
    与证据流建模在动态运维场景中的能力差异。
\end{enumerate}


% =========================================================
\subsection{相关工作与研究差异}
\label{sec:related_work}
% =========================================================

\textbf{检索增强生成与结构化上下文组织。}
经典 RAG 将外部非参数知识库与生成模型结合，
通过检索相关文档为生成阶段提供可追溯的外部知识
\cite{lewis2020rag}。
随后研究逐渐从独立短文本检索扩展到更加结构化的上下文组织。
RAPTOR 递归地对文本块进行聚类与摘要，
构建多层次树状索引，
使查询能够在不同语义粒度上获取信息
\cite{sarthi2024raptor}；
GraphRAG 则从语料中构建实体知识图和社区摘要，
利用图结构支撑跨文档的信息聚合
\cite{edge2024graphrag}。
这类方法表明，检索单元之间的组织结构会直接影响最终上下文质量。
但其结构主要由文档层级或语料实体关系预先构建，
而 TEF-RAG 的关系图在查询时刻由当前可见运维证据动态形成，
并显式表示状态更新、规程约束、诊断支持及结果验证等有方向的关系。

\textbf{时间感知 RAG 与动态知识检索。}
动态知识环境下，
仅依据语义相似度容易召回时间上不适用的事实。
MRAG 将问题中的主题内容和时间约束分离，
通过语义--时间联合排序处理 time-sensitive question answering
\cite{zhang2025mrag}。
TA-RAG 面向需要跨时间观察实体变化的 diachronic questions，
进一步强调证据在目标时间区间上的连续覆盖
\cite{lau2025tarag}。
本文与这些工作的共同点是显式建模时间条件，
但问题设定更强调\emph{知识可见性}而不仅是时间相关性：
一条较早发生的记录如果在查询截止时刻之后才进入系统，
仍不能被当前查询使用。
此外，TEF-RAG 不仅判断单条证据是否处于目标时间范围，
还要求所选多条证据能够通过有向关系形成完整的任务相关证据流。

\textbf{多跳证据检索与集合级选择。}
复杂问答通常需要联合多个互补证据。
Multi-Hop Dense Retrieval 通过前一跳结果条件化后续检索，
在无需预定义语料图结构的情况下收集多跳证据
\cite{xiong2021mdr}。
这类方法的核心是逐步扩展检索路径；
TEF-RAG 则不把目标限定为单一线性路径，
而是先在候选证据之间建立查询条件化的 typed directed graph，
再对固定大小的证据集合进行联合排序。
因此，本文的重点不是``下一跳文档是哪一条''，
而是``一个候选集合能否同时覆盖必要角色，
并在时间方向和关系结构上组成可用于当前决策的完整证据流''。
这一区别也决定了本文采用集合级 Complete/FlowComplete 指标，
而不是仅以单条文档相关性或逐跳命中率衡量检索质量。

\textbf{面向机器人/具身运维的 RAG 与电池时序辅助决策。}
检索结果在本文中最终服务于具身运维任务规划，而非停留在问答层。
SayComply 将操作要求、环境条件与机器人能力组织为层级知识，
利用检索结果约束现场机器人任务满足操作合规要求
\cite{ginting2025saycomply}；
RAGLRO 进一步在电力开关柜操作场景中利用操作手册与安全协议
生成机器人控制指令，验证了 RAG 与专业设备操作规划结合的可行性
\cite{wang2026raglro}。
在电池与储能相关任务中，LiBrain 强调由专业时序模型承担预测、
RAG 提供领域知识、LLM 进行综合推理的分工
\cite{li2026librain}；
TimeSeries2Report 则探索将多变量电池时序数据转换为结构化语义报告，
再供语言模型进行状态管理与决策
\cite{yang2026timeseries2report}。
这些工作与本文的应用目标一致，即将外部知识与动态状态最终转化为可执行任务。
本文的研究差异在于：不直接把生成器作为主要创新对象，
而是集中解决生成之前的证据选择问题，
使标准化工单与有序具身指令建立在查询时刻可见、
关系一致且流程完整的证据上下文之上。


